\scanpage{77}
\mspage{34}
\uline{15}. Introduction des proespaces :
$V_{i_0\ldots i_n},\ V_J^{\deleted{\unsure{*}}\unsure{L}}$,

On désigne par $V_i^{\unsure{\deleted{*}|*}}$ le
\inserted{pro}voisinage tubulaire (au sens des proespaces)
de $X_i^*=X_i\setminus\bigcup_{j<i}X_j$ dans $X$.
Ainsi, \deleted{\unsure{on}} on a
\[
(51)\qquad X_i^*\longrightarrow V_i^{\unsure{\deleted{*}|*}}
\qquad\text{équivalence d'homotopie de proespaces.}
\]
\deleted{Pour un drapeau}
Si $\unsure{\forall i,j|(i,j)}$ $i<j$, on pose
\[
(52)\qquad\deleted{X_{ij}}\inserted{V_{ij}}=V_i\cap V_j
\]
\margin{Pourquoi
\diagram{\begin{tikzcd}[column sep=small]
V_{ij}^{\unsure{b}}=V_i\cap X_j^*\arrow[r,hook]&{}
\end{tikzcd}}
\unsure{est une|est-ce une} équivalence d'homotopie ?}

\margin{NB $V_{ij}$ a un sens \deleted{\unsure{si}} si $i\neq j$
mais si $i$ et $j$ ne sont pas comparables, on a $V_i\cap V_j=\varnothing$,
\inserted{(du moins si les ouverts de $X$ sont normaux, p.ex. $X$ métrisable)}
si $i=j$, on a $V_{ii}=V_i$}

de sorte qu'on a
\diagram{\begin{tikzcd}[column sep=large,row sep=large]
&V_{ij}\arrow[r,"b_{ij}"]\arrow[d,leftrightarrow,"s_{ij}"]&V_j\\
(X_i\quad\deleted{\longleftarrow}\arrow[r,"{\substack{(w)\\\sim}}"]&)V_i&
\end{tikzcd}}
On a :
\[
(53)\qquad X=\varinjlim(V_i,V_{ij},\deleted{\unsure{\}}}\ b_{ij},s_{ij})
\]
$X$ est obtenu par \uline{recollement} des $V_i$ suivant les $V_{ij}$.

On définit aussi, plus généralement, les
\[
(54)\qquad V_{i_0\ldots i_n}=V_{i_0}\cap\ldots\cap V_{i_n}
\]
pour un drapeau $i_0<i_1<\ldots<i_n$
\margin{\uline{NB} L'expression a un sens \unsure{même} si
l'ens. $\{i_0,\ldots,i_n\}$ ne forme pas un drapeau,
mais il est vide si \unsure{ce} n'est pas un drapeau}
\mspage{76}
